\begin{enumerate}
\item Sí, en $\mathbb{R}$.
\item No es solución en ningún intervalo abierto no vacío: la igualdad sólo se cumple en $x=0$.
\item No en $\mathbb{R}$; sí en $(-\infty,0)$ y $(0,\infty)$. No es derivable en $0$.
\item Sí: define una única función derivable en $\mathbb{R}$, con $y'=(1+\cos x)/(1+3y^2)$.
\item La curva completa no es una gráfica. Cada rama $y=\pm\sqrt{\ln(1+x)}$ es solución en $(0,\infty)$; el punto $(0,0)$ no puede pertenecer a una solución de la EDO.
\item Sí, como función de $x$ en $\mathbb{R}$: $y'=2t/(1+3t^2)$.
\item Sí, en $\mathbb{R}$: la gráfica es $y=x^2+x$. En $t=0$ el cociente paramétrico es indeterminado, pero $y'(0)=1$.
\end{enumerate}
